% !TEX program = lualatex
% Transparencias de Fundamentos de las Matemáticas II
% Clase 07: Factorización en primos I

\documentclass[aspectratio=169,11pt]{beamer}

\usepackage{fontspec}
\usepackage[spanish,es-nodecimaldot]{babel}
\usepackage{amsmath,amssymb,mathtools}
\usepackage{booktabs}
\usepackage{csquotes}
\usepackage{xcolor}

\usetheme{Madrid}
\usecolortheme{default}
\setbeamertemplate{navigation symbols}{}
\setbeamertemplate{footline}[frame number]

\definecolor{FdMBlue}{RGB}{20,55,100}
\definecolor{FdMRed}{RGB}{130,30,30}
\definecolor{FdMGreen}{RGB}{25,100,60}

\setbeamercolor{structure}{fg=FdMBlue}
\setbeamercolor{title}{fg=white,bg=FdMBlue}
\setbeamercolor{frametitle}{fg=white,bg=FdMBlue}
\setbeamercolor{block title}{fg=white,bg=FdMBlue}
\setbeamercolor{block body}{bg=blue!3}
\setbeamercolor{alerted text}{fg=FdMRed}

\setmainfont{Latin Modern Roman}
\setsansfont{Latin Modern Sans}
\setmonofont{Latin Modern Mono}

\newcommand{\N}{\mathbb{N}}
\newcommand{\Z}{\mathbb{Z}}
\newcommand{\Q}{\mathbb{Q}}
\newcommand{\R}{\mathbb{R}}
\newcommand{\C}{\mathbb{C}}
\newcommand{\Pow}{\mathcal{P}}
\newcommand{\id}{\operatorname{id}}
\newcommand{\mcd}{\operatorname{mcd}}
\newcommand{\mcm}{\operatorname{mcm}}
\newcommand{\im}{\operatorname{Im}}

\title[FdM II -- Clase 07]{Factorización en primos I}
\subtitle{Existencia de la factorización}
\author{Antonio Falcó Montesinos}
\institute{Universidad CEU Cardenal Herrera}
\date{Fundamentos de las Matemáticas II}

\begin{document}

\begin{frame}
  \titlepage
\end{frame}

\begin{frame}{Objetivos de la clase}
\begin{itemize}
  \item Enunciar la existencia de factorización en primos.
  \item Usar inducción o buen orden en la prueba.
  \item Distinguir primos y compuestos.
\end{itemize}
\end{frame}

\begin{frame}{Mapa de la clase}
\tableofcontents
\end{frame}

\section{Factorización}

\begin{frame}{Factorización}
\begin{block}{Objetivo}
\begin{itemize}
  \item Todo entero \(n>1\) puede escribirse como producto de números primos.
  \item La escritura puede tener un solo factor si \(n\) ya es primo.
  \item La existencia es la primera mitad del Teorema Fundamental de la Aritmética.
\end{itemize}
\end{block}
\end{frame}

\begin{frame}{Factorización}
\begin{block}{Idea de prueba}
\begin{itemize}
  \item Si \(n\) es primo, ya está factorizado.
  \item Si \(n\) es compuesto, entonces \(n=ab\) con \(1<a<n\) y \(1<b<n\).
  \item Se aplica inducción a \(a\) y \(b\).
\end{itemize}
\end{block}
\end{frame}

\begin{frame}{Factorización}
\begin{block}{Ejemplo}
\begin{itemize}
  \item \(84=2\cdot 42=2\cdot 2\cdot 21=2^2\cdot 3\cdot 7\).
  \item El proceso termina porque los factores disminuyen.
  \item La descomposición final contiene solo primos.
\end{itemize}
\end{block}
\end{frame}

\section{Herramientas}

\begin{frame}{Herramientas}
\begin{block}{Inducción fuerte}
\begin{itemize}
  \item La existencia de factorización es un ejemplo natural de inducción fuerte.
  \item Para factorizar \(n\), puede ser necesario conocer factorizaciones de varios menores que \(n\).
  \item El buen orden ofrece una prueba alternativa por mínimo contraejemplo.
\end{itemize}
\end{block}
\end{frame}

\begin{frame}{Herramientas}
\begin{block}{Papel del número 1}
\begin{itemize}
  \item El \(1\) no es primo.
  \item Si lo fuera, la factorización no sería única.
  \item Por ejemplo, \(6=2\cdot 3=1\cdot 2\cdot 3=1^2\cdot 2\cdot 3\).
\end{itemize}
\end{block}
\end{frame}

\section{Profundización}

\begin{frame}{Profundización}
\begin{block}{Mínimo contraejemplo}
\begin{itemize}
  \item Supongamos que existe un entero \(>1\) sin factorización.
  \item Por buen orden, existe uno mínimo, \(m\).
  \item Si \(m\) es compuesto, sus factores menores sí factorizan, contradicción.
\end{itemize}
\end{block}
\end{frame}

\begin{frame}{Profundización}
\begin{block}{Práctica}
\begin{itemize}
  \item Descomponer \(126\).
  \item Descomponer \(360\).
  \item Explicar en cada paso por qué el proceso termina.
\end{itemize}
\end{block}
\end{frame}

\begin{frame}{Profundización}
\begin{block}{Control}
\begin{itemize}
  \item ¿Dónde se usa que los factores son menores que \(n\)?
  \item ¿Por qué no basta encontrar una descomposición cualquiera?
  \item ¿Qué falta todavía para el teorema completo?
\end{itemize}
\end{block}
\end{frame}


\section{Detalle y práctica}

\begin{frame}{Detalle y práctica}
\begin{block}{Primos y compuestos}
\begin{itemize}
  \item Un entero \(p>1\) es primo si sus únicos divisores positivos son \(1\) y \(p\).
  \item Un entero \(n>1\) es compuesto si admite un divisor positivo distinto de \(1\) y \(n\).
  \item Para probar que \(n\) es compuesto basta exhibir una factorización no trivial.
\end{itemize}
\end{block}
\end{frame}

\begin{frame}{Detalle y práctica}
\begin{block}{Actividad breve}
\begin{itemize}
  \item Clasificar \(29\), \(51\), \(91\) y \(97\) como primos o compuestos.
  \item Justificar por qué basta probar divisores primos \(\leq \sqrt{n}\).
  \item Explicar por qué \(1\) no se considera primo.
\end{itemize}
\end{block}
\end{frame}

\begin{frame}{Cierre}
\begin{itemize}
  \item La existencia se prueba naturalmente por inducción fuerte.
  \item El número \(1\) queda excluido de los primos.
  \item Aún falta demostrar la unicidad.
\end{itemize}
\end{frame}

\begin{frame}{Trabajo recomendado}
\begin{itemize}
  \item Releer las secciones correspondientes del manual.
  \item Identificar definiciones, símbolos nuevos y enunciados principales.
  \item Preparar dudas y ejemplos para el seminario asociado.
\end{itemize}
\end{frame}

\end{document}
